\scanpage{25}
\mspage{24}

En effet, $\widetilde{F\vee\unsure{G|C}}=\widetilde F\cup\widetilde G$, et $\widetilde F$ et $\widetilde G$ y sont fermés, \deleted{dans $M$} ce qui prouve que
\[
 \unsure{\mathrm{Multomb}|\mathrm{Multiomb}}(F\unsure{\cup|\vee}G)=\mathrm{Multomb}(F)\cup\mathrm{Mulomb}(G),
\]
et que $\mathrm{Mulomb}(F)$ et $\mathrm{Mulomb}(G)$ y sont fermés, donc \unsure{(*)} et la première relation (**). La deuxième relation \inserted{(**)} provient de $=$ de $\widetilde{F\cap G}=\widetilde F\cap\widetilde G$.

\uline{NB} Il n'est pas vrai en général que $\mathrm{Mulomb}(F) \mathrel{\tikz[baseline=-.55ex,x=1em,y=1em]{\draw[line width=.065em,line cap=round,line join=round] (0,.40)--(.48,0)--(0,-.35) (.51,.40)--(.03,.02)--(.49,-.35);}} \mathrm{Mulomb}(G)$ implique $F \mathrel{\tikz[baseline=-.55ex,x=1em,y=1em]{\draw[line width=.065em,line cap=round,line join=round] (0,.40)--(.48,0)--(0,-.35) (.51,.40)--(.03,.02)--(.49,-.35);}} G$, même si $F$ et $G$ sont réduits aux lieux distincts $x$ et $y$ --- auquel cas $\mathrm{Multomb}(X)=\{\{\unsure{\lambda|x}\}\}$, $\mathrm{Mulomb}(Y)=\{\{y\}\}$ sont compatibles et $=$ disjointes, mais $x$ et $y$ ne le sont pas forcément. Mais il \unsure{sera|ser.} vrai \unsure{par contre qu'on a équivalence|pour d'autres quasi-équivalences|pour d'autre quasi-équivalence}, dans le cas où $\unsure{\mathcal F|\mathscr F}$ est \unsure{un|une} \unsure{atelier} \unsure{ensembliste|ens.embliste}.
% slanted
\margin{\unsure{vérifier|vérifies} plus tard}

\uline{Corollaire 2} Si $F$ et $G$ sont disjoints, alors $\unsure{\mathrm{Mul.mb}|\mathrm{Mulomb}}(F)$ et $\mathrm{Mulomb}(G)$ aussi.

Cela provient formellement du corollaire 1 et de la formule
\[
 \mathrm{Mulomb}(\varnothing_{\mathcal F})=\varnothing\quad\text{(figure vide dans $M$)}
 \tag*{(1.40)}
\]

\uline{Proposition} \unsure{Soient|Soit} $\unsure{F|\mathbb F},G\in\mathcal F$. Alors $F\ll G$
\[
 \Longleftrightarrow\quad\mathrm{Mulomb}(F)\ll\mathrm{Mulomb}(G).
\]
